% Living Showcase: Abbildungen, Plots und Side-by-side (Palette Slot 5).
\StartChapter[ch:figures]{Abbildungen und Layout}

\SubsectionBar[sec:fig-helpers]{Bild-Makros}
\ZSFindex{ZSFfig}
\ZSFindex{ZSFfigside}

\ZSFfig[width=0.34]{example-image-a}{Ein Bild, optional skaliert\ZSFTitleTag{ZSFfig}}%
  {width und height sind Schlüssel wie überall; ohne sie greifen die zentralen Vorbelegungen.}

\ZSFfigside[width=0.27]{example-image-a}{Bild links, Text rechts\ZSFTitleTag{ZSFfigside}}{%
  Der erste Wert bestimmt den Bildanteil, der Rest gehört dem Text.}

\ZSFfig{example-image-a,example-image-b,example-image-c}%
  {Dasselbe Makro, mehrere Pfade\ZSFTitleTag{ZSFfig}}%
  {Eine Pfadliste ergibt eine Reihe; die Breite folgt der Anzahl.}

\SubsectionBar[sec:plot]{Plot ohne externe Datei}
\ZSFindex{figbox}

\begin{figbox}[Diagramm direkt im Dokument\ZSFTitleTag{figbox + pgfplots}]
  \centering
  \begin{tikzpicture}
    \begin{axis}[
      width=\linewidth,
      height=\ZSFfigMaxHeight,
      axis lines=middle,
      xmin=-2.2, xmax=2.2,
      ymin=-0.3, ymax=4.5,
      samples=40,
      ticks=none
    ]
      \addplot[color=\chaptercolor, very thick, domain=-2:2]{x^2};
      % Relative Achsenkoordinaten statt absoluter: die Beschriftung bleibt
      % dadurch garantiert innerhalb der Box und trifft die freie Fläche
      % zwischen den Parabelästen, statt am Rand abgeschnitten zu werden.
      \node[\chaptercolor, font=\ZSFfontDiagramLabel, anchor=center]
        at (rel axis cs:0.68,0.72) {ZSFfontDiagramLabel};
    \end{axis}
  \end{tikzpicture}
  \ZSFfigcaption{Knotenbeschriftung über ZSFfontDiagramLabel statt lokalem
    \textbackslash tiny; pgfplots kommt aus dem Opt-in-Modul 12\_plots.}
\end{figbox}

\begin{defbox}[Derselbe Regler, mit Kopfzeile\ZSFTitleTag{split}][split=0.34]
  \centering Formel
  \tcblower
  \ZSFkeyword*{split} ist ein Regler wie jeder andere und gilt deshalb auch
  auf einer Box mit Titel — vorher brauchte es dafür eine Verschachtelung.
\end{defbox}

\begin{splitbox}[split=0.34, sidebyside align=center]
  \centering linker Block
  \tcblower
  \ZSFkeyword{splitbox} ordnet zwei beliebige Inhaltsblöcke ohne sichtbaren
  Rahmen nebeneinander an; \ZSFkeyword*{split} legt den linken Anteil fest.
\end{splitbox}
